\documentclass[11pt]{article}

\usepackage[a4paper,margin=1in]{geometry}
\usepackage{amsmath,amssymb,amsthm,mathtools}
\usepackage{booktabs}
\usepackage{enumitem}
\usepackage{microtype}
\usepackage{xcolor}
\usepackage{hyperref}

\hypersetup{
    colorlinks=true,
    linkcolor=blue!60!black,
    citecolor=blue!60!black,
    urlcolor=blue!60!black,
    pdftitle={Signatures Before Roots},
    pdfauthor={Vinicius Santos}
}

\newtheorem{definition}{Definition}[section]
\newtheorem{theorem}[definition]{Theorem}
\newtheorem{corollary}[definition]{Corollary}
\newtheorem{proposition}[definition]{Proposition}
\newtheorem{lemma}[definition]{Lemma}
\newtheorem{observation}[definition]{Observation}
\newtheorem{question}[definition]{Question}
\newtheorem{remark}[definition]{Remark}

\newcommand{\B}{\mathsf{B}}
\newcommand{\Tr}{\operatorname{Tr}}
\newcommand{\Nm}{\operatorname{N}}
\newcommand{\Sp}{\operatorname{Sp}}
\newcommand{\Disc}{\operatorname{Disc}}
\newcommand{\Fix}{\operatorname{Fix}}
\newcommand{\Lift}{\operatorname{Lift}}
\newcommand{\RootLift}{\operatorname{RootLift}}
\newcommand{\LenLift}{\operatorname{LenLift}}
\newcommand{\sig}{\sigma}
\newcommand{\ph}{\varphi}
\newcommand{\ps}{\psi}
\newcommand{\Z}{\mathbb{Z}}
\newcommand{\Q}{\mathbb{Q}}
\newcommand{\R}{\mathbb{R}}
\newcommand{\C}{\mathbb{C}}

\title{Signatures Before Roots\\
\large Shadow Lifts, Norm Forms, and Radical Spreads}
\author{Vinicius Santos\\
\small Exploratory draft}
\date{June 2026}

\begin{document}
\maketitle

\begin{abstract}
The companion manuscripts \emph{A Golden Shadow Algebra}, \emph{Quotients Without Division}, and \emph{Shadow Trigonometry} suggest a recurring principle: operations that appear primitive in ordinary notation may become projections, relations, constraints, or lift problems in a $\ph$-native algebraic presentation.  Addition becomes trace projection.  Quotients split into unit actions, rate localizations, relations, approximations, and factor projections.  Trigonometry begins from norm-one phase objects and only later acquires angles, normalized sine, normalized cosine, and $\pi$.

This note applies the same strategy to roots.  Its central identity is the quadratic shadow formula
\[
    \Sp^2=\Tr^2-4\Nm.
\]
For a pair $\langle r_1\mid r_2\rangle$ with trace $S=r_1+r_2$ and norm $P=r_1r_2$, the hidden spread $D=r_1-r_2$ satisfies
\[
    D^2=S^2-4P.
\]
Thus the classical radical $\sqrt{S^2-4P}$ is not treated as a primitive scalar function.  It is the recovery of the anti-fixed spread component from visible invariant data.

The sign and square-class of the discriminant play different roles.  Over a real embedding, the sign of $\Delta=S^2-4P$ determines the signature of the lift: elliptic/circular when $\Delta<0$, degenerate/null when $\Delta=0$, and hyperbolic/split when $\Delta>0$.  Algebraically, the square-class of $\Delta$ determines whether the spread already lives in the current scalar layer or requires an extension.  This distinction explains the double life of the golden ratio.  As a circular phase trace, $(S,P)=(\ph,1)$ and $\Delta=\ph-3<0$.  As the real golden pair $\langle\ph\mid\ps\rangle$, $(S,P)=(1,-1)$ and $\Delta=5>0$.  The same shadow-lift mechanism produces the pentagonal circle in one signature and the golden hyperbola in another.
\end{abstract}

\tableofcontents

\section{Introduction: do not start with the radical}

The previous notes in this sequence began from a simple methodological rule: do not force a familiar operation into the wrong layer.  In \emph{A Golden Shadow Algebra}, the scalar grammar generated by
\[
    \B(x,y)=xy-y
\]
from the golden leaf $\ph$ was shown to be multiplication/predecessor-native but not addition-native.  Addition appeared only after passing to a pair layer:
\[
    x+y=\Tr\langle x\mid y\rangle.
\]
In \emph{Quotients Without Division}, division split into several roles: unit action, rate localization, quotient relation, reciprocal approximation, normalization constraint, and factor projection.  In \emph{Shadow Trigonometry}, trigonometry was not begun with angles, sine, cosine, or tangent.  It began with norm-one phase objects and their unnormalized trace and spread projections:
\[
    T(U)=U+U^{-1},
    \qquad
    S(U)=U-U^{-1}.
\]

The present note applies the same discipline to roots.  Classically, one writes expressions such as
\[
    \sqrt{x},
    \qquad
    \sqrt{a^2+b^2},
    \qquad
    \sqrt{S^2-4P}.
\]
These notations hide several different tasks.  Sometimes a root is a length extraction.  Sometimes it is a discriminant square root.  Sometimes it is a request to pass from symmetric coefficients to individual roots.  Sometimes it is an embedding choice: real, complex, split, or degenerate.

The guiding question is:
\begin{quote}
What if a square root is not a primitive scalar function, but the shadow of a missing component?
\end{quote}

For a quadratic pair $\langle r_1\mid r_2\rangle$, the visible invariant data are trace and norm:
\[
    S=r_1+r_2,
    \qquad
    P=r_1r_2.
\]
The hidden anti-fixed data is the spread:
\[
    D=r_1-r_2.
\]
The fundamental identity is
\[
    D^2=S^2-4P.
\]
Thus the radical $\sqrt{S^2-4P}$ is not the first object.  It appears when one tries to recover the spread from its square.

The main thesis is:
\begin{quote}
Signature first, spread second, root later.
\end{quote}

More precisely, the real sign of $S^2-4P$ tells us whether the lift is circular, null, or hyperbolic under a chosen real embedding.  The algebraic square-class of $S^2-4P$ tells us whether the spread is already present in the current scalar layer or belongs to a further extension.  Roots therefore become lift requests, not primitive total functions.

\section{The native hierarchy ladder}

This fourth note uses every layer introduced so far.  It is useful to make the ladder explicit.

\begin{center}
\begin{tabular}{p{0.27\linewidth}p{0.28\linewidth}p{0.34\linewidth}}
\toprule
Layer & Role & Native mechanism \\
\midrule
Scalar layer & constants, unit actions, multiplication & $\B(x,y)=xy-y$, leaf $\ph$, constants $\Z[\ph]$ \\
Trace layer & addition & pair formation and $\Tr\langle x\mid y\rangle=x+y$ \\
Rate layer & fixed-denominator division & localization classes and diagonal trace sections \\
Factor layer & derivative-like quotients & factor projection from divisibility relations \\
Phase layer & trigonometry & norm-one phase, torsion, trace/spread projections \\
Signature layer & roots and embeddings & spread lifts from $\Delta=S^2-4P$ \\
\bottomrule
\end{tabular}
\end{center}

The present paper occupies the last row.  It does not replace the earlier layers.  It depends on them.  In particular, recovering individual root coordinates from trace and spread often uses fixed averaging by $2$, which belongs to the rate/localization layer of \emph{Quotients Without Division}.  The root mechanism itself is the spread lift; coordinate extraction is a later layer.

\section{Trace, norm, and spread}

We begin with the same pair language used throughout the sequence.

\begin{definition}[Quadratic shadow pair]
Let $A$ be a commutative ring and let
\[
    R=\langle r_1\mid r_2\rangle
\]
be a formal pair with slot-swap involution
\[
    \sig\langle r_1\mid r_2\rangle=\langle r_2\mid r_1\rangle.
\]
Define
\[
    \Tr(R)=r_1+r_2,
\]
\[
    \Nm(R)=r_1r_2,
\]
and
\[
    \Sp(R)=r_1-r_2.
\]
\end{definition}

The trace and norm are invariant under the slot-swap involution.  The spread is anti-invariant:
\[
    \sig(\Sp)=-\Sp.
\]
This is why spread is the hidden component in a symmetric presentation: trace and norm live downstairs, while spread records the sign-sensitive distinction between the two slots.

\begin{proposition}[Spread-discriminant identity]\label{prop:spread-discriminant}
For a quadratic shadow pair $R=\langle r_1\mid r_2\rangle$,
\[
    \Sp(R)^2=\Tr(R)^2-4\Nm(R).
\]
\end{proposition}

\begin{proof}
Directly,
\[
    (r_1-r_2)^2
    =r_1^2-2r_1r_2+r_2^2
    =(r_1+r_2)^2-4r_1r_2.
\]
This is precisely
\[
    \Sp(R)^2=\Tr(R)^2-4\Nm(R).
\]
\end{proof}

This identity is the central theorem package of the note.  It is also the familiar discriminant identity for a monic quadratic.  If
\[
    S=r_1+r_2,
    \qquad
    P=r_1r_2,
\]
then $r_1,r_2$ are the slots of the Vieta lift
\[
    X^2-SX+P=0,
\]
and
\[
    \Delta=S^2-4P=(r_1-r_2)^2.
\]
The novelty is not the formula itself.  The formula is classical.  The point is its role in the shadow hierarchy: the radical is the spread asking to be lifted.

\section{Roots as spread lifts}

In classical notation, one often writes
\[
    \sqrt{S^2-4P}.
\]
In the shadow formulation we instead define a lift relation.

\begin{definition}[Spread lift]
Let $S,P$ be visible invariant data.  A spread lift of $(S,P)$ is an element $D$ satisfying
\[
    D^2=S^2-4P.
\]
We write
\[
    \Lift_{\Sp}(S,P;D)
\]
for this relation.
\end{definition}

The spread lift does not automatically produce ordered root slots.  If $2$ is invertible, or if one works in the rate/localization layer where fixed averaging by $2$ is available, then
\[
    r_1=\frac{S+D}{2},
    \qquad
    r_2=\frac{S-D}{2}.
\]
But this coordinate reconstruction is a separate operation.

\begin{remark}[The role of $2$]
The spread relation $D^2=S^2-4P$ makes sense without solving for $r_1$ and $r_2$ as scalar coordinates.  Recovering the ordered slots from $S$ and $D$ uses fixed averaging by $2$.  In characteristic $2$, this split degenerates.  Over integral coefficient layers such as $\Z[\ph]$, the division by $2$ is not scalar-native; it belongs to the rate/localization layer described in \emph{Quotients Without Division}.  Thus the root mechanism is spread lift, while coordinate extraction is an averaging operation.
\end{remark}

This distinction is important.  The shadow-native object is not the formula
\[
    \frac{S\pm \sqrt{S^2-4P}}{2}
\]
viewed as a scalar expression.  It is the pair-lift problem:
\[
    \text{given }(S,P),\text{ find slots }\langle r_1\mid r_2\rangle
\]
with trace $S$ and norm $P$.

\section{Real signature and algebraic square-class}

There are two different questions one can ask about
\[
    \Delta=S^2-4P.
\]
They are related but not identical.

\subsection{Real signature}

After choosing a real embedding, the sign of $\Delta$ classifies the real geometry of the lift.

\begin{center}
\begin{tabular}{lll}
\toprule
Real sign of $\Delta$ & Lift type & Geometry \\
\midrule
$\Delta<0$ & elliptic / circular & conjugate phase, no real spread \\
$\Delta=0$ & degenerate / parabolic & double slot, null boundary \\
$\Delta>0$ & hyperbolic / split & two real slots, real spread \\
\bottomrule
\end{tabular}
\end{center}

For example, the equation
\[
    X^2-SX+P=0
\]
has two real roots when $\Delta>0$, one repeated real root when $\Delta=0$, and a conjugate pair over a quadratic imaginary extension when $\Delta<0$.

\subsection{Square-class}

Over a field $K$, the finer algebraic question is not only whether $\Delta$ is positive or negative under one embedding, but whether it is a square in $K$.  The relevant invariant is the square-class
\[
    \Delta\in K^\times/(K^\times)^2.
\]
If $\Delta$ is a square in $K$, then the spread already exists in $K$.  If not, the spread lives in an extension
\[
    K(\sqrt{\Delta}).
\]

Over $\R$, sign and square-class collapse into the familiar dichotomy: positive elements are squares and negative elements are not.  Over the field $\Q(\ph)$, the square-class lives in $K^\times/(K^\times)^2$.  Over the ring $\Z[\ph]$, one must separately ask whether an element is literally a square in the ring, a stronger integrality condition not captured by field square-class alone.  An element may be positive in one real embedding, negative in another, or positive but still not a square in the field or in the ring.  Thus this paper keeps two words separate:
\begin{itemize}[leftmargin=2em]
    \item \emph{signature} refers to real geometric sign behavior;
    \item \emph{square-class} refers to algebraic solvability of the spread lift in a field layer; ring-level squareness is a further integrality question.
\end{itemize}

This distinction prevents a common confusion.  A real picture may suggest a circle or hyperbola, while the algebraic layer may still require an extension to realize the spread.

\section{Embedding as signature selection}

The tangent discussion in \emph{Shadow Trigonometry} required a complex embedding, or equivalently a skew unit $i$ satisfying
\[
    i^2=-1,
    \qquad
    i^\sig=-i.
\]
In that manuscript this appeared as a necessary hand-off to a coordinate layer.  The spread-discriminant viewpoint explains why the hand-off is not arbitrary.

Consider the circular base pair with
\[
    S=0,
    \qquad
    P=1.
\]
Then
\[
    \Delta=S^2-4P=-4.
\]
The spread $D$ satisfies
\[
    D^2=-4.
\]
After fixed averaging by $2$, define
\[
    i=\frac{D}{2}.
\]
Then
\[
    i^2=-1.
\]
Thus the imaginary unit is the normalized spread generator of the elliptic signature.  The complex embedding is not a random import.  It is the coordinate realization of the negative-discriminant spread lift.

\begin{proposition}[The skew unit as elliptic spread]\label{prop:skew-unit-spread}
Let $D$ be a spread lift for $(S,P)=(0,1)$ over a layer where fixed averaging by $2$ is available.  Then $i=D/2$ satisfies $i^2=-1$ and is anti-fixed under the shadow involution.
\end{proposition}

\begin{proof}
Since $D^2=-4$, fixed averaging gives
\[
    \left(\frac{D}{2}\right)^2=-1.
\]
The spread is anti-fixed because slot-swap sends $r_1-r_2$ to $r_2-r_1$.  Hence $i^\sig=-i$.
\end{proof}

This reframes a possible weakness as structure.  Tangent was not smuggling in $i$; tangent was asking to work in the elliptic signature where the spread has square $-4$.  The embedding is selected by the discriminant.

\section{The two golden anchors}

The discriminant viewpoint explains the double life of the golden ratio.  The same symbol $\ph$ appears in two different quadratic shadow configurations.  The trace and norm knobs are different in the two cases.

\subsection{Golden circular phase}

In \emph{Shadow Trigonometry}, the golden phase was defined by
\[
    U+U^{-1}=\ph.
\]
This is the Vieta data
\[
    S=\ph,
    \qquad
    P=1.
\]
Thus
\[
    \Delta_{\mathrm{phase}}
    =S^2-4P
    =\ph^2-4
    =(\ph+1)-4
    =\ph-3<0.
\]
So the golden phase lift is elliptic/circular.  It satisfies
\[
    U^2-\ph U+1=0,
\]
and the previous paper showed
\[
    U^5=-1,
    \qquad
    U^{10}=1.
\]
The golden ratio is therefore a visible trace of a hidden pentagonal circular phase.

\subsection{Golden scalar pair}

The real golden scalar pair is
\[
    \Phi=\langle\ph\mid\ps\rangle,
    \qquad
    \ps=1-\ph.
\]
Here the Vieta data are
\[
    S=\ph+\ps=1,
    \qquad
    P=\ph\ps=-1.
\]
Therefore
\[
    \Delta_{\mathrm{scalar}}
    =1^2-4(-1)=5>0.
\]
This lift is hyperbolic/split.  Its spread is real:
\[
    \Sp(\Phi)=\ph-\ps=\sqrt5.
\]

\begin{center}
\begin{tabular}{lllll}
\toprule
Case & $S$ & $P$ & $\Delta=S^2-4P$ & Signature \\
\midrule
Golden phase & $\ph$ & $1$ & $\ph-3<0$ & circular / elliptic \\
Golden scalar pair & $1$ & $-1$ & $5>0$ & hyperbolic / split \\
Circular base pair & $0$ & $1$ & $-4<0$ & elliptic, gives $i$ \\
\bottomrule
\end{tabular}
\end{center}

The same mechanism produces different geometries.  The difference is not only trace and not only norm.  The signature is controlled by both:
\[
    \Delta=S^2-4P.
\]

\section{Norms come in signatures}

We now shift from the discriminant of a pair to the norm form of an algebra.  Again, the key point is signature.

\subsection{The golden scalar norm}

Let
\[
    z=a+b\ph\in\Z[\ph].
\]
Under the real quadratic involution
\[
    \ph^\sig=\ps=1-\ph,
\]
we have
\[
    z^\sig=a+b\ps.
\]
The norm is
\[
    \Nm_\ph(z)=zz^\sig=(a+b\ph)(a+b\ps).
\]
Using
\[
    \ph+\ps=1,
    \qquad
    \ph\ps=-1,
\]
we obtain
\[
    \Nm_\ph(a+b\ph)=a^2+ab-b^2.
\]
This form is indefinite.  Its level sets
\[
    a^2+ab-b^2=c
\]
are hyperbolas, not circles.

The null directions are obtained from
\[
    a+b\ph=0
    \qquad\text{or}\qquad
    a+b\ps=0,
\]
which are the two real shadow directions of the quadratic field.  Thus the golden scalar norm is naturally Lorentzian/hyperbolic.

\subsection{Positive-definite skew norms}

By contrast, if one adjoins a skew unit $\eta$ satisfying
\[
    \eta^\sig=-\eta,
    \qquad
    \eta^2=-d,
    \qquad
    d>0,
\]
then
\[
    \Nm(a+b\eta)=(a+b\eta)(a-b\eta)=a^2+d b^2.
\]
This norm is positive definite.  The Gaussian case corresponds to $d=1$ and $\eta=i$:
\[
    \Nm(a+bi)=a^2+b^2.
\]

This skew unit is a genuine adjunction.  It is not already present in $\Z[\ph]$.  In the language of Proposition~\ref{prop:skew-unit-spread}, it is the coordinate realization of an elliptic spread lift.

\section{Pythagoras as positive-definite norm}

The Pythagorean theorem belongs naturally to the positive-definite norm case.  In the Gaussian algebra,
\[
    z=a+bi,
    \qquad
    z^\sig=a-bi,
\]
we have
\[
    \Nm(z)=a^2+b^2.
\]
Thus a right triangle relation
\[
    a^2+b^2=c^2
\]
is the norm-square relation
\[
    \Nm(a+bi)=c^2.
\]
This suggests a native formulation.

\begin{definition}[Length lift]
A length lift of an element $z$ in a positive-definite norm layer is an element $c$ satisfying
\[
    c^2=\Nm(z).
\]
We write
\[
    \LenLift(z;c)
\]
for this relation.
\end{definition}

The important point is that Pythagoras does not begin with the scalar expression
\[
    c=\sqrt{a^2+b^2}.
\]
It begins with the norm identity
\[
    \Nm(a+bi)=a^2+b^2.
\]
The scalar length $c$ is a later lift.

In the golden scalar norm, however,
\[
    \Nm_\ph(a+b\ph)=a^2+ab-b^2
\]
is indefinite.  There are three regions under a real embedding:

\begin{center}
\begin{tabular}{lll}
\toprule
Region & Condition & Root-lift behavior \\
\midrule
Spacelike & $\Nm_\ph(z)>0$ & real length lift may exist \\
Null & $\Nm_\ph(z)=0$ & boundary directions \\
Timelike & $\Nm_\ph(z)<0$ & no real length lift in the base real layer \\
\bottomrule
\end{tabular}
\end{center}

So Pythagoras is not the universal norm story.  It is the positive-definite signature of a broader norm-first geometry.

\section{The Fibonacci unit hyperbola}

The golden norm has its own native geometry.  Its unit level sets are Pell-type hyperbolas:
\[
    a^2+ab-b^2=\pm1.
\]
The integer points are controlled by the units of $\Z[\ph]$.

\begin{lemma}[Powers of the golden unit]\label{lem:phi-powers-fib}
For $k\ge1$,
\[
    \ph^k=F_{k-1}+F_k\ph,
\]
where $F_0=0,F_1=1$ are the Fibonacci numbers.  Negative powers have the coordinate form
\[
    \ph^{-k}=(-1)^k\bigl(F_{k+1}-F_k\ph\bigr)
    \qquad (k\ge1).
\]
\end{lemma}

\begin{proof}
The formula for positive powers follows by induction.  For $k=1$ this says $\ph=F_0+F_1\ph$.  If
\[
    \ph^k=F_{k-1}+F_k\ph,
\]
then
\[
    \ph^{k+1}=F_{k-1}\ph+F_k\ph^2
    =F_{k-1}\ph+F_k(\ph+1)
    =F_k+(F_{k-1}+F_k)\ph
    =F_k+F_{k+1}\ph.
\]
For negative powers, use $\ph^{-1}=\ph-1$ and the identity $\Nm(\ph^k)=(-1)^k$; equivalently, conjugation gives $\ps^k=F_{k-1}+F_k\ps$, while $\ps=-\ph^{-1}$.  Solving for $\ph^{-k}$ yields the displayed formula.
\end{proof}

Since
\[
    \Nm(\ph)=\ph\ps=-1,
\]
we get
\[
    \Nm(\ph^k)=(-1)^k.
\]
Using the positive-power part of Lemma~\ref{lem:phi-powers-fib}, this becomes the Fibonacci hyperbola identity
\[
    F_{k-1}^2+F_{k-1}F_k-F_k^2=(-1)^k
    \qquad(k\ge1).
\]

\begin{theorem}[Fibonacci points on the golden unit hyperbola]\label{thm:fibonacci-hyperbola}
The integral points on the norm levels
\[
    a^2+ab-b^2=\pm1
\]
are precisely the coordinate pairs of the units
\[
    \pm\ph^k\qquad(k\in\Z).
\]
For $k\ge1$, the positive powers have coordinates
\[
    \ph^k \leftrightarrow (F_{k-1},F_k),
\]
and the negative powers have coordinates
\[
    \ph^{-k} \leftrightarrow \bigl((-1)^kF_{k+1},(-1)^{k+1}F_k\bigr).
\]
Multiplying by $-1$ gives the remaining sign-symmetric points.
\end{theorem}

\begin{proof}
The displayed coordinate formulas follow from Lemma~\ref{lem:phi-powers-fib}.  Their norms are $\pm1$ by multiplicativity of the norm.  Conversely, an integral element $a+b\ph$ has norm $\pm1$ if and only if it is a unit of $\Z[\ph]$.  In this quadratic order, $\ph$ is a fundamental unit and the unit group is
\[
    \{\pm\ph^k:k\in\Z\}.
\]
Thus every integral point on the two unit norm levels is one of the displayed golden-unit points.
\end{proof}

This is the golden-native analogue of the Pythagorean-triple story.  It is not Euclidean; it is hyperbolic.  The Fibonacci numbers appear as lattice coordinates on the unit hyperbola.

\section{Gaussian triples as contrast}

The classical Pythagorean triples belong to the imaginary quadratic positive-definite case.  In $\Z[i]$,
\[
    \Nm(a+bi)=a^2+b^2.
\]
A Pythagorean triple satisfies
\[
    \Nm(a+bi)=c^2.
\]
When
\[
    a+bi=(m+ni)^2,
\]
we obtain
\[
    a=m^2-n^2,
    \qquad
    b=2mn,
    \qquad
    c=m^2+n^2.
\]
This is classical.  Its role here is comparative: Gaussian Pythagoras is the positive-definite norm-square case, while the golden scalar geometry is the indefinite Fibonacci-hyperbola case.

The comparison is useful precisely because it shows that ``norm before root'' is not one geometry.  It is a framework whose geometry is determined by signature.

\section{Vieta lifts and polynomial roots}

The quadratic formula is often written as
\[
    r=\frac{S\pm\sqrt{S^2-4P}}{2}.
\]
In the shadow formulation, this is decomposed into layers:

\begin{enumerate}[label=(\roman*)]
    \item visible symmetric data $(S,P)$;
    \item spread lift $D^2=S^2-4P$;
    \item rate/averaging by $2$ to recover coordinates;
    \item slot choice, deciding which root is $r_1$ and which is $r_2$.
\end{enumerate}

The root pair itself is the Vieta lift
\[
    \langle r_1\mid r_2\rangle
\]
with
\[
    r_1+r_2=S,
    \qquad
    r_1r_2=P.
\]
This is the classical symmetric-function viewpoint in shadow language.  Coefficients live downstairs as symmetric data.  Individual roots live in a lifted object.

\begin{remark}[No novelty claim for Vieta]
Nothing in this section is meant as a new theorem about quadratic equations.  The value is organizational: the square root in the quadratic formula is identified with the spread lift, and the division by $2$ is assigned to the rate/localization layer.  This assignment is what keeps the native hierarchy explicit.
\end{remark}

\section{Higher degree: discriminants, not spreads}

The spread mechanism is special to degree $2$.  For a quadratic, the hidden anti-symmetric direction is one-dimensional:
\[
    r_1-r_2.
\]
The symmetric group $S_2$ has a one-dimensional sign representation, so a single spread captures the anti-fixed component.

For degree $n\ge3$, the situation changes.  A root tuple
\[
    \langle r_1\mid r_2\mid\cdots\mid r_n\rangle
\]
has elementary symmetric invariants
\[
    e_1,e_2,\dots,e_n.
\]
The discriminant remains a visible symmetric invariant:
\[
    \Disc=\prod_{i<j}(r_i-r_j)^2.
\]
But there is no single spread that recovers all hidden root data.  The individual differences $r_i-r_j$ are acted on by the Galois group, and lifting from symmetric data to individual roots is the full splitting-field problem.

Thus the degree-$2$ slogan
\[
    \sqrt{\Delta}=\text{spread lift}
\]
should not be naively generalized.  In higher degree, the discriminant is still a shadow of anti-symmetric root data, but recovering the slots is a Galois problem.  In particular, solvability by radicals is a special property, not a general guarantee.

\section{The floor: where lifting stops}

The phrase ``the radical is a shadow asking to be lifted'' can sound like an infinite regress.  The series has a floor: the finite scalar grammar of \emph{A Golden Shadow Algebra}.  At that base layer, constants are exactly
\[
    \Z[\ph],
\]
and scalar terms are polynomial, division-free, and root-free.  Every later operation is explicitly assigned to a higher layer: trace, rate localization, factor projection, phase lift, or spread lift.

This is the sense in which the framework is falsifiable.  A proposed radical may or may not be available in the current layer.  If it is not, one must say which extension, lift, or localization has been added.  Roots are not silently totalized.

\section{Comparison with classical root notation}

\begin{center}
\begin{tabular}{p{0.32\linewidth}p{0.52\linewidth}}
\toprule
Classical notation & Shadow-native interpretation \\
\midrule
$\sqrt{x}$ & root-lift relation $r^2=x$ \\
$\sqrt{S^2-4P}$ & spread lift $D^2=S^2-4P$ \\
Quadratic formula & Vieta lift + spread lift + rate averaging by $2$ \\
$\sqrt{a^2+b^2}$ & length lift from positive-definite norm \\
$2\cos\theta$ & trace of norm-one phase \\
$2i\sin\theta$ & spread of norm-one phase \\
$i=\sqrt{-1}$ & normalized spread generator of elliptic signature \\
Golden units & integer points on the Fibonacci hyperbola \\
\bottomrule
\end{tabular}
\end{center}

The recurring theme is the same as in the previous papers: ordinary notation often compresses several structural steps into one symbol.  The shadow-native approach expands the symbol into its native layers.

\section{Open questions}

\begin{question}[Square-classes in $\Q(\ph)$]
Given $x\in\Q(\ph)$, when is $x$ a square in the field $\Q(\ph)$?  If $x\in\Z[\ph]$, when is it already a square in the ring $\Z[\ph]$?  How do these field square-class and ring-level squareness questions interact with the two real embeddings of $\Q(\ph)$?
\end{question}

\begin{question}[Golden norm geometry]
Can the geometry of the indefinite form
\[
    a^2+ab-b^2
\]
be developed as a useful $\ph$-native analogue of Euclidean norm geometry, with its own notions of null direction, orthogonality, and geodesic flow?
\end{question}

\begin{question}[Spread lifts and phase lifts]
How should one formally relate spread lifts $D^2=S^2-4P$ to phase lifts $U^2-SU+P=0$ in a typed algebraic system?  Are these two presentations term-equivalent after adjoining the rate layer for division by $2$?
\end{question}

\begin{question}[Higher-degree anti-symmetric data]
For cubic and quartic shadow tuples, which anti-symmetric or representation-theoretic data best replaces the quadratic spread?  How much of classical Galois theory can be reformulated as a hierarchy of shadow lifts?
\end{question}

\begin{question}[Calculus implications]
If distance and roots are lift relations rather than primitive scalar functions, how should gradients, curvature, and normalization be formulated in a $\ph$-native calculus?
\end{question}

\section{Conclusion: the radical is a shadow asking to be lifted}

Classically one writes
\[
    r=\sqrt{x}.
\]
This note argues that the more native question is:
\[
    \text{what object has }x\text{ as the square of its hidden component?}
\]
For a quadratic shadow pair,
\[
    \Sp^2=\Tr^2-4\Nm.
\]
The discriminant is not merely an expression under a radical.  It is the square of the spread.

The sign of the discriminant selects a real signature: circular, null, or hyperbolic.  Its square-class selects whether the spread lives in the current algebraic layer or demands an extension.  The golden ratio exhibits both sides of the story.  As a phase trace, it gives a negative discriminant and pentagonal circular geometry.  As the real scalar pair $\langle\ph\mid\ps\rangle$, it gives discriminant $5$ and hyperbolic golden geometry.

Pythagoras then finds its place as a positive-definite norm case, not the root of the whole theory.  The golden scalar norm is indefinite, and its native unit geometry is the Fibonacci hyperbola.

Thus the fourth step of the program is:
\begin{quote}
Roots are not primitive scalar functions.  They are lift requests from visible invariant data to hidden shadow components.
\end{quote}

Final slogan:
\[
    \boxed{\text{the radical is a shadow asking to be lifted.}}
\]

\begin{thebibliography}{9}

\bibitem{SantosGoldenShadow}
V.~Santos.
\newblock \emph{A Golden Shadow Algebra: From a $\varphi$-Leaf Scalar Core to Conjugate-Pair Arithmetic}.
\newblock Exploratory manuscript, 2026.

\bibitem{SantosQuotients}
V.~Santos.
\newblock \emph{Quotients Without Division: Unit Actions, Relations, and Factor Projections in a $\varphi$-Native Algebra}.
\newblock Exploratory manuscript, 2026.

\bibitem{SantosShadowTrig}
V.~Santos.
\newblock \emph{Shadow Trigonometry: Norm-One Orbits, Golden Traces, and Cyclotomic Phase}.
\newblock Exploratory manuscript, 2026.

\bibitem{Odrzywolek2026EML}
A.~Odrzywo\l{}ek.
\newblock \emph{All elementary functions from a single binary operator}.
\newblock arXiv:2603.21852, 2026.

\bibitem{Marcus1977}
D.~A. Marcus.
\newblock \emph{Number Fields}.
\newblock Springer-Verlag, 1977.

\bibitem{Washington1997}
L.~C. Washington.
\newblock \emph{Introduction to Cyclotomic Fields}.
\newblock 2nd edition, Springer, 1997.

\bibitem{IrelandRosen1990}
K.~Ireland and M.~Rosen.
\newblock \emph{A Classical Introduction to Modern Number Theory}.
\newblock 2nd edition, Springer, 1990.

\bibitem{StewartGalois}
I.~Stewart.
\newblock \emph{Galois Theory}.
\newblock 4th edition, Chapman and Hall/CRC, 2015.

\end{thebibliography}

\end{document}
